% Pillar-2 (ZVF) results opener
\section{The ZVF Diagnostic: Results}
\label{sec:p2-results}
We evaluate ZVF along three axes: its association with training failure, its
temporal behavior, and its coupling to nuisance variables. Pooling per-experiment
means across the benchmark (using per-experiment rather than per-step rows, since
per-step ZVF is autocorrelated), mean ZVF correlates with catastrophic collapse
at Spearman $\rho \approx 0.56$--$0.62$ but with the final held-out accuracy at
only $\rho \approx 0.27$. The gap is the crux of the pillar: ZVF flags when a run
has stopped producing within-group signal, yet a low ZVF does not guarantee a
large held-out gain, because a group can be information-rich and still fail to
generalize.

The subsections that follow develop this picture. We first formalize ZVF and its
measurement pipeline, then show it is a temporally sticky signal: across 60
measured runs spanning nine variance-mitigation methods and the group-size sweep,
late-phase ZVF ($0.87$--$0.97$ for GRPO) far exceeds early-phase ZVF
($0.04$--$0.13$), with lag-1 autocorrelation $\approx 0.94$, consistent with a
signal-starvation trajectory rather than a property of the task. We connect ZVF to
the gradient through the utilization identity $\mathrm{GU} = 1 - \mathrm{ZVF}$, and
audit its coupling to reward sparsity, group size, and baseline accuracy so that
its correlations are read as descriptive rather than causal. The final subsection
presents an external frontier-model cross-examination that isolates ZVF's core
defect---it assigns the same value to a mastered prompt and an unsolvable one---and
proposes magnitude- and sign-aware replacements (a pairwise-contrast density and a
learning-adjusted rollout quality score).
